\subsection{\texorpdfstring{Closed ideals of L \(^{1}\) (G)}{Closed ideals of L \^{}\{1\} (G)}}

The Fourier cotransform on the Banach algebra \(L^{1}(G)\) is identified with the Gelfand transform of \(L^{1}(G)\) (II, p.~209). With this identification, recall that if I is an ideal of \(L^{1}(G)\) we denote by V(I) the closed set in \(\widehat{G}\) of characters \(\chi \in \widehat{G}\) such that, for every function \(f \in I\) , the Fourier cotransform of f vanishes at \(\chi\) (cf.~I, p.~30). For every subset M of \(\widehat{G}\) , we denote by \(\Upsilon(M)\) the closed ideal consisting of the \(f \in L^{1}(G)\) such that \(\overline{\mathcal{F}}_{\mathrm{G}}(f)\) vanishes on M (I, p.~30).

By prop. 2 of II, p.~219, the Banach algebra L \(^{1}\) (G) is regular. By § 5 of I, p.~88, we therefore deduce the following properties of the Fourier transform and cotransform:

\begin{enumerate}
\def\labelenumi{\arabic{enumi})}
\item
  If F is a closed subset of \(\widehat{\mathbf{G}}\) and K is a compact subset of \(\widehat{\mathbf{G}}\) such that \(\mathrm{F} \cap \mathrm{K} = \varnothing\) , there exists a function \(f \in \mathrm{L}^{1}(\mathrm{G})\) such that \(\mathcal{F}_{\mathrm{G}}(f)\) is equal to 0 on F and to 1 on K (I, p.~88, prop. 1; for this fact and the following ones, one passes from the Fourier cotransform to the Fourier transform by means of formula (8) of II, p.~207).
\item
  Let M be a closed subset of \(\widehat{G}\) . The set of ideals I of \(L^{1}(G)\) such that \(V(I) = M\) has \(\Upsilon(M)\) as its greatest element and, as its smallest element, the set of \(f \in \mathrm{L}^{1}(\mathrm{G})\) whose Fourier cotransform has compact support disjoint from M (I, p.~91, prop. 4).
\item
  Let I be an ideal of \(L^{1}(G)\) , and \(g\colon\widehat{G}\to C\) a continuous function. We assume that for every \(\chi\in\widehat{G}\) , there exists a function \(f_{\chi}\in I\) such that g is equal to \(\mathcal{F}_{\mathrm{G}}(f_{\chi})\) in a neighborhood of \(\chi\). We assume in addition that there exists a function \(f_{\infty}\in I\) such that g is equal to \(\mathcal{F}_{\mathrm{G}}(f_{\infty})\) in the complement of a compact subset of \(\widehat{G}\) , this last condition being always satisfied if G is discrete. Then there exists a function \(f\in I\) such that \(g=\mathcal{F}_{\mathrm{G}}(f)\) (I, p.~91, cor. 2).
\end{enumerate}

Lemma 1. --- The space of functions in \(\mathrm{L}^{1}(\mathrm{G})\) whose Fourier transform has compact support is dense in \(\mathrm{L}^{1}(\mathrm{G})\) .

Since \(\mathcal{K}(\widehat{\mathrm{G}})\) is dense in \(\mathrm{L}^2 (\widehat{\mathrm{G}})\) and since the Fourier transform of \(\mathrm{L}^2 (\mathrm{G})\) is an isometry on \(\mathrm{L}^2 (\widehat{\mathrm{G}})\) (Theorem 1 of II, p.~215), the subspace V of \(\mathrm{L}^2 (\mathrm{G})\) consisting of the \(f\in \mathrm{L}^2 (\mathrm{G})\) such that \(\mathcal{F}_{\mathrm{G}}(f)\in \mathcal{K}(\widehat{\mathrm{G}})\) is dense in \(\mathrm{L}^2 (\mathrm{G})\) .

Let \(g \in \mathrm{L}^{1}(\mathrm{G})\) . There exists \(g_{1}, g_{2} \in \mathrm{L}^{2}(\mathrm{G})\) such that \(g = g_{1}g_{2}\) (one can, for example, take \(g_{1} = |g|^{1/2}\) , and \(g_{2}(x) = 0\) if \(g(x) = 0\) , \(g_{2}(x) = g(x)/g_{1}(x)\) otherwise). From the above, we therefore deduce that g is the limit of a sequence of functions of the form \(h_{1}h_{2}\) , where \(h_{1}\) and \(h_{2}\) belong to V. However, \(\mathcal{F}_{\mathrm{G}}(h_{1}h_{2}) = \mathcal{F}_{\mathrm{G}}(h_{1}) * \mathcal{F}_{\mathrm{G}}(h_{2})\) (II, p.~223, prop. 14), and \(\mathcal{F}_{\mathrm{G}}(h_{1}) * \mathcal{F}_{\mathrm{G}}(h_{2})\) belongs to \(\mathcal{K}(\mathrm{G})\) . The lemma follows.

PROPOSITION 2. --- Let I be a closed ideal of L \(^{1}\) (G), and let f ∈ L \(^{1}\) (G). If \(\overline{\mathcal{F}}_{\text{G}}(f)\) vanishes on a neighborhood of V(I), then f belongs to I.

Let \(\varepsilon > 0\) . There exists a function \(g \in \mathrm{L}^1(\mathrm{G})\) such that \(\|f - f * g\|_1 < \varepsilon\) (prop. 8 of II, p.~211 (iv)). Let \(h \in \mathrm{L}^1(\mathrm{G})\) such that the support of \(\overline{\mathcal{F}}_{\mathrm{G}}(h)\) is compact and \(\|f\|_1 \|g - h\|_1 < \varepsilon\) (lemma 1). We have

\[
\| f - f * h \| _ {1} \leqslant \| f - f * g \| _ {1} + \| f * (g - h) \| _ {1} <   2 \varepsilon .
\]

By the hypothesis on \(f\) , the function \(\overline{\mathcal{F}}_{\mathrm{G}}(f*h)=\overline{\mathcal{F}}_{\mathrm{G}}(f)\overline{\mathcal{F}}_{\mathrm{G}}(h)\) is compactly supported and disjoint from V(I), which implies that \(f*h\in\mathrm{I}\) (remark 2 above). Since \(\varepsilon\) is arbitrarily small, we have \(f\in\overline{\mathrm{I}}=\mathrm{I}\) .

THEOREM 1. --- Let I be a closed ideal of L \(^{1}\) (G) distinct from L \(^{1}\) (G). There exists a character \(\widehat{x} \in \widehat{\mathrm{G}}\) such that \(\mathcal{F}_{\mathrm{G}}(f)(\widehat{x}) = 0\) for every \(f \in \mathrm{I}\) .

Since the algebra L \(^{1}\) (G) is regular (prop. 2 of II, p.~219) and semisimple (cor. of prop. 22 of I, p.~126), and since the set of functions whose Fourier cotransform has compact support is dense in

\(L^{1}(G)\) (lemma 1), cor. 1 of I, p.~92 shows that the ideal I is contained in a regular maximal ideal of \(L^{1}(G)\) , that is, in the kernel of a character \(\widehat{y}\) of \(L^{1}(G)\) (th. 2 of I, p.~30); one can then take \(\widehat{x} = \widehat{y}^{-1}\) (cf.~formula (8) of II, p.~207).

COROLLARY 1 (Wiener's tauberian theorem)

Let \(f \in \mathrm{L}^1(\mathrm{G})\) . If the Fourier transform of \(f\) does not vanish, the functions \(f * \varepsilon_x \colon g \mapsto f(gx^{-1})\) , where \(x\) ranges over \(\mathrm{G}\) , form a total set in \(\mathrm{L}^1(\mathrm{G})\) (EVT, I, p.~12, def. 1).

Let V be the closed vector subspace of \(L^{1}(G)\) generated by the \(f*\varepsilon_{x}\) . By INT, VIII, §4, cor. of prop. 20, the space V is a closed ideal of \(L^{1}(G)\) . By th. 1 we have \(V = L^{1}(G)\) .

DEFINITION 1. --- Let g be a complex function on G and let Φ be a filter on G. We say that g is slowly oscillating along Φ if, for every ε \textgreater{} 0, there exists a set M ∈ Φ and a neighborhood V of e in G such that

\[
 x \in \mathrm{M} \quad \text {and} \quad y \in \mathrm{V} \quad \Longrightarrow \quad | g (x y) - g (x) | \leqslant \varepsilon .
\]

COROLLARY 2. --- Let \(\Phi\) be a filter on G invariant under translation. Let \(f \in \mathrm{L}^{1}(\mathrm{G})\) such that the Fourier transform of \(f\) does not vanish and such that \(\int_{\mathrm{G}} f(x) dx = 1\) . Let \(g \in \mathrm{L}^{\infty}(\mathrm{G})\) . We assume that \(f * g\) has a finite limit \(\alpha\) along \(\Phi\) .

\begin{enumerate}
\def\labelenumi{\alph{enumi})}
\item
  For every function \(h \in \mathrm{L}^{1}(\mathrm{G})\) such that \(\int_{\mathrm{G}} h(x) dx = 1\) , the limit of \(h * g\) along \(\Phi\) is equal to \(\alpha\) ;
\item
  Suppose moreover that \(g\) is slowly oscillating along \(\Phi\) . Then \(g\) tends to \(\alpha\) along \(\Phi\) .
\end{enumerate}

By replacing \(g\) by \(g - \alpha\) , one reduces to the case where \(\alpha = 0\) . Let I be the set of functions \(h \in \mathrm{L}^1(\mathrm{G})\) such that \(h*g\) tends to 0 along \(\Phi\) . The set I is a translation-invariant vector subspace of \(\mathrm{L}^1(\mathrm{G})\) . It is a closed vector subspace. Indeed, let \(h \in \bar{\mathrm{I}}\) . For every function \(h_0 \in \mathrm{L}^1(\mathrm{G})\) and every \(x \in \mathrm{G}\) , one has

\[
\begin{array}{c} | (h * g) (x) | \leqslant | ((h - h _ {0}) * g) (x) | + | (h _ {0} * g) (x) | \\ \leqslant \| h - h _ {0} \| _ {1} \| g \| _ {\infty} + | (h _ {0} * g) (x) |. \end{array}
\]

For every \(\varepsilon > 0\) , there exists \(h_0 \in \mathrm{I}\) such that \(\|h - h_0\|_1 \|g\|_\infty < \varepsilon\) . Let \(\mathrm{M} \in \Phi\) such that \(|(h_0 * g)(x)| < \varepsilon\) for all \(x \in \mathrm{M}\) . We then have \(|(h * g)(x)| < 2\varepsilon\) for all \(x \in \mathrm{M}\) , hence \(h * g\) converges to 0 along \(\Phi\) . This shows that \(h \in \mathrm{I}\) .

The space I is therefore a closed ideal of \(\mathrm{L}^{1}(\mathrm{G})\) . We have \(f \in I\) by hypothesis, therefore \(\mathrm{I} = \mathrm{L}^{1}(\mathrm{G})\) by th. 1, since the Fourier transform of f does not vanish. This implies a).

Assume now the hypotheses of \(b\) ). Let \(\varepsilon > 0\) . Since \(g\) is slowly oscillating along \(\Phi\) , there exists \(\mathrm{M} \in \Phi\) and a compact neighborhood \(\mathrm{V}\) of \(e\) such that

\[
 x \in \mathrm{M} \quad \text { and } \quad y \in \mathrm{V} \Longrightarrow | g (y ^ {- 1} x) - g (x) | \leqslant \varepsilon .
\]

Let \(\varphi\) be the characteristic function of V and \(\mu = \int \varphi(x)dx\) . For every \(x \in G\) , one has

\[
\frac {1}{\mu} (\varphi * g) (x) = \frac {1}{\mu} \int_ {\mathrm{V}} g (y ^ {- 1} x) d y = g (x) + \frac {1}{\mu} \int_ {\mathrm{V}} (g (y ^ {- 1} x) - g (x)) d y.
\]

Therefore for all \(x \in \mathrm{M}\) , one has

\[
\left| \frac {1}{\mu} (\varphi * g) (x) - g (x) \right| \leqslant \varepsilon .
\]

Since, by \(a\) ), the limit of \(\varphi * g\) along \(\Phi\) is zero, we have \(\limsup_{\Phi} |g| \leqslant \varepsilon\) . Since \(\varepsilon\) is arbitrary, we conclude that the limit of \(g\) along \(\Phi\) is zero, which proves \(b\) ).

Lemma 2. --- Let K be a compact subset of G. For every \(\eta > 0\) , there exists a function \(j \in \mathrm{L}^{1}(\mathrm{G})\) such that:

\begin{enumerate}
\def\labelenumi{\alph{enumi})}
\item
  \(\|j\|_{1}\leqslant\sqrt{2};\)
\item
  the function \(\mathcal{F}_{\mathrm{G}}(j)\) is equal to 1 in a neighborhood of the identity element of \(\widehat{\mathrm{G}}\) ;
\item
  for all \(x \in \mathbf{K}\) , one has \(\| j - j * \varepsilon_x\|_1 \leqslant \eta\) .
\end{enumerate}

The set U \(_{1}\) of elements \(\widehat{x} \in \widehat{\mathrm{G}}\) such that

\[
| \langle \widehat {x}, x \rangle - 1 | \leqslant \frac {\eta}{4}
\]

 for all \(x \in K\) is a neighborhood of e in \(\hat{G}\) . Let \(U \subset U_{1}\) be an open, symmetric neighborhood integrable with respect to Haar measure \(m = d\hat{x}\) of \(\hat{G}\) dual to the measure dx. Let \(V \subset U\) be a compact symmetric neighborhood of e such that \(m(V) \geqslant \frac{1}{2}m(U)\) . Denote by \(\varphi_{U}\) (resp. \(\varphi_{V}\) ) the characteristic function of U (resp. of V). Since \(\varphi_{U}\) belongs to \(L^{2}(G)\) , there exists \(u \in L^{2}(G)\) such that \(\varphi_{U} = \mathcal{F}_{G}(u)\) (th. 1 of II, p.~215). Similarly, there exists a function \(v \in L^{2}(G)\) such that \(\varphi_{V} = \mathcal{F}_{G}(v)\) . We shall show that the function \(j = \frac{1}{m(V)}uv\) satisfies the required properties. We have \(j \in L^{1}(G)\) .

\begin{enumerate}
\def\labelenumi{\alph{enumi})}
\tightlist
\item
  By the Plancherel theorem and the condition \(m(\mathrm{V}) \geqslant \frac{1}{2}m(\mathrm{U})\) , one has
\end{enumerate}

\[
\| j \| _ {1} \leqslant \frac {\| u \| _ {2} \| v \| _ {2}}{m (\mathrm{V})} = \frac {\| \mathcal {F} _ {\mathrm{G}} (u) \| _ {2} \| \mathcal {F} _ {\mathrm{G}} (v) \| _ {2}}{m (\mathrm{V})} = \frac {\sqrt {m (\mathrm{U}) m (\mathrm{V})}}{m (\mathrm{V})} \leqslant \sqrt {2}.
\]

\begin{enumerate}
\def\labelenumi{\alph{enumi})}
\setcounter{enumi}{1}
\tightlist
\item
  There exists a neighborhood W of e in \(\widehat{G}\) such that WV ⊂ U (TG, II, p.~31, prop. 4). For every \(\widehat{x} \in W\) , one has \(\widehat{x}V \subset U\) , and prop. 14 of II, p.~223 implies
\end{enumerate}

\[
\begin{array}{l} \mathcal {F} _ {\mathrm{G}} (j) (\widehat {x}) = \frac {1}{m (\mathrm{V})} (\mathcal {F} _ {\mathrm{G}} (u) * \mathcal {F} _ {\mathrm{G}} (v)) (\widehat {x}) \\ \qquad = \frac {1}{m (\mathrm{V})} \int_ {\widehat {\mathrm{G}}} \varphi_ {\mathrm{U}} (\widehat {y}) \varphi_ {\mathrm{V}} (\widehat {y} ^ {- 1} \widehat {x}) d m (\widehat {y}) \\ \qquad = \frac {m (\mathrm{U} \cap \widehat {x} \mathrm{V} ^ {- 1})}{m (\mathrm{V})} = \frac {m (\widehat {x} \mathrm{V})}{m (\mathrm{V})} = 1 \end{array}
\]

since V is symmetric.

\begin{enumerate}
\def\labelenumi{\alph{enumi})}
\setcounter{enumi}{2}
\tightlist
\item
  If \(x \in \mathrm{K}\) , one has
\end{enumerate}

\[
\| u - u * \varepsilon_ {x} \| _ {2} ^ {2} = \int_ {\widehat {\mathrm{G}}} \big | \mathcal {F} _ {\mathrm{G}} (u) (\widehat {x}) \big (1 - \overline {{\langle x , \widehat {x} \rangle}} \big) \big | ^ {2} d m (\widehat {x}) \leqslant m (\mathrm{U}) \Big (\frac {\eta}{4} \Big) ^ {2}
\]

since \(\mathrm{U}\subset \mathrm{U}_1\) , and likewise \(\| v - v*\varepsilon_x\| _2^2\leqslant m(\mathrm{V})\big(\frac{\eta}{4}\big)^2\) . Therefore

\[
\begin{array}{l} \| j - j * \varepsilon_ {x} \| _ {1} = \frac {1}{m (\mathrm{V})} \| u (v - v * \varepsilon_ {x}) + (v * \varepsilon_ {x}) (u - u * \varepsilon_ {x}) \| _ {1} \\ \leqslant \frac {\eta}{4 m (\mathrm{V})} (\| u \| _ {2} \sqrt {m (\mathrm{V})} + \| v \| _ {2} \sqrt {m (\mathrm{U})}) \\ = \frac {\eta \sqrt {m (\mathrm{U}) m (\mathrm{V})}}{2 m (\mathrm{V})} <   \eta . \end{array}
\]

PROPOSITION 3. --- The algebra L \(^{1}\) (G) satisfies the Ditkin condition (I, p.~92, def. 2).

Let \(\chi\) be a character of \(\mathrm{L}^1 (\mathrm{G})\) . We distinguish two cases according as \(\chi\) is zero or not. If \(\chi\) is zero, it must be verified that for every function \(f\in \mathrm{L}^{1}(\mathrm{G})\) there exists a sequence \((f_n)_{n\geqslant 1}\) in \(\mathrm{L}^1 (\mathrm{G})\) such that \(\mathcal{F}(f_n)\) vanishes outside a compact subset of \(\widehat{\mathbf{G}}\) and such that \(f_{n}*f\) tends to \(f\) in \(\mathrm{L}^1 (\mathrm{G})\) . The existence of such a sequence follows from Lemma 1 above and from Prop. 8 of II, p.~211.

Suppose now that \(\chi\) is nonzero, hence \(\chi \in \mathsf{X}(\mathrm{L}^1 (\mathrm{G})) = \widehat{\mathrm{G}}\) (prop. 1 of II, p.~202). Let \(f\in \mathrm{L}^{1}(\mathrm{G})\) such that \(\mathcal{G}_{\mathrm{L}^1 (\mathrm{G})}(f)(\chi) = \overline{\mathcal{F}} (f)(\chi) = 0\) . It remains to prove the existence of a sequence \((f_n)_{n\geqslant 1}\) in

\(L^{1}(G)\) such that \(f * f_{n}\) converges to f in \(L^{1}(G)\) and such that \(\overline{\mathcal{F}}(f_{n})\) vanishes in a neighborhood of \(\chi\) . One may assume that \(\|f\|_{1}=1\) . By translation in \(\widehat{G}\) , one reduces to the case where \(\chi=e\) .

Let \(K_{n}\) be a compact subset of G such that

\[
\int_ {\mathrm{G} - \mathrm{K} _ {n}} | f (x) | d x \leqslant \frac {1}{n}.
\]

Let \(u_{n} \in L^{1}(G)\) be a function \(\geqslant 0\) such that \(\|u_{n}\|_{1} = 1\) and

\[
\| f - f * u _ {n} \| _ {1} \leqslant \frac {1}{n}
\]

(cf. prop. 8 of II, p. 211, (iii)). By Lemma 2, there exists a function \(j_n\) in \(\mathrm{L}^1(\mathrm{G})\) such that \(\| j_n\| _1\leqslant \sqrt{2}\) whose Fourier cotransform equals 1 in a neighborhood of \(e\) and moreover such that \(\| j_{n} - j_{n}*\varepsilon_{x}\|_{1}\leqslant n^{-1}\) for all \(x\in \mathrm{K}_n\) . We set

\[
f _ {n} = u _ {n} - j _ {n} * u _ {n}.
\]

We shall show that the sequence \((f_{n})_{n\geqslant1}\) has the required properties. First of all, we have

\[
\mathcal {F} (f _ {n}) = \mathcal {F} (u _ {n}) - \mathcal {F} (j _ {n}) \mathcal {F} (u _ {n}) = (1 - \mathcal {F} (j _ {n})) \mathcal {F} (u _ {n})
\]

so the Fourier transform of \(f_{n}\) vanishes in a neighborhood of \(\chi = e\) . On the other hand,

\[
\| f * f _ {n} - f \| _ {1} \leqslant \| f * u _ {n} - f \| _ {1} + \| f * j _ {n} \| _ {1} \| u _ {n} \| _ {1} \leqslant \frac {1}{n} + \| f * j _ {n} \| _ {1}.
\]

Now, for almost every \(y \in G\) , one has

\[
(f * j _ {n}) (y) = \int_ {\mathrm{G}} f (x) j _ {n} (x ^ {- 1} y) d x = \int_ {\mathrm{G}} f (x) (j _ {n} (x ^ {- 1} y) - j _ {n} (y)) d x
\]

since, by hypothesis, we have \(\mathcal{F}(f)(e) = \int_{\mathrm{G}} f(x) dx = 0\) . Hence

\[
\begin{array}{l} \| f * j _ {n} \| _ {1} \leqslant \int_ {\mathrm{G}} | f (x) | \| j _ {n} * \varepsilon_ {x} - j _ {n} \| _ {1} d x \\ \qquad = \int_ {\mathrm{K} _ {n}} | f (x) | \| j _ {n} * \varepsilon_ {x} - j _ {n} \| _ {1} d x \\ \qquad + \int_ {\mathrm{G-K} _ {n}} | f (x) | \| j _ {n} * \varepsilon_ {x} - j _ {n} \| _ {1} d x \\ \qquad \leqslant \frac {1}{n} \int_ {\mathrm{K} _ {n}} | f (x) | d x + 4 \int_ {\mathrm{G-K} _ {n}} | f (x) | d x \leqslant \frac {5}{n}. \end{array}
\]

Finally, \(\|f * f_{n} - f\|_{1} \leqslant 6n^{-1}\) and therefore \(f * f_{n}\) converges to \(f\) in \(L^{1}(G)\) , as desired.

Applying prop. 5, we then obtain the following result:

THEOREM 2. --- Let I be a closed ideal of L \(^{1}\) (G) such that the boundary of V(I) contains no nonempty perfect set. Then I is the set of functions \(f \in \text{L}^{1}(\text{G})\) such that \(\mathcal{F}(f)\) vanishes on V(I).

For any closed ideal of \(L^{1}(G)\) , the conclusion of Thm. 2 is in general false (cf.~Exerc. 12 of II, p.~314). More precisely, one can show that if G is noncompact, there exists a closed ideal of \(L^{1}(G)\) which is not self-adjoint (see, for example, W. RUDIN, Fourier analysis on groups, Interscience tracts in pure and applied mathematics, theorem 7.7.1.)

COROLLARY. --- If a closed ideal I of L \(^{1}\) (G) is contained in a single regular maximal ideal, then I is itself regular maximal.

